% !PRINTED 642
\weberpage{642}
%<<CONT
friedigt sein kann, wenn alle ihre Co\"efficienten verschwinden. Dies f\"uhrt zu den Gleichungen
\[
x_{\frac12\mu}=0,\ x_{\frac12\mu-1}=x_{-\frac12\mu+1},\ \dots,\ x_1=x_{-1},
\]
und es ist also jede ganze Zahl $\omega$ des K\"orpers $H_m$ in der Form enthalten:
\[
\omega = x_0+x_1(r+r^{-1})+x_2(r^2+r^{-2})+\cdots
+x_{\frac12\mu-1}(r^{\frac12\mu-1}+r^{-\frac12\mu+1}).
\]

Statt der Basis (1) kann man auch, wenn man
\[
\varrho = r+r^{-1}
\]
setzt, die Potenzen von $\varrho$:
\weq{3}{1,\ \varrho,\ \varrho^2,\ \dots,\ \varrho^{\frac12\mu-1}}
als Basis der ganzen Zahlen von $H_m$ w\"ahlen. Denn die Gr\"ossen (3) lassen sich ganzzahlig durch die (1) ausdr\"ucken und umgekehrt.

Die Zahl $\varrho$ ist die Wurzel einer irreduciblen Gleichung vom Grade $\frac12\mu$, die wir mit $\psi(\varrho)=0$ bezeichnen wollen. Ist $f(r)=0$ die Gleichung f\"ur $r$ (\S.\,169), so besteht die identische Relation
\weq{4}{f(x) = x^{\frac12\mu}\,\psi\Bigl(x+\frac1x\Bigr),}
woraus durch Bildung der Ableitung
\weq{5}{f'(r) = r^{\frac12\mu}\,\psi'(\varrho)\,(1-r^{-2}).}

Hieraus k\"onnen wir die Grundzahl $\Delta_1$ des K\"orpers $H_m$ bilden, die, da $H_m$ nur reelle Zahlen enth\"alt, positiv sein muss. Diese Grundzahl ist, wenn wir die Norm in Bezug auf $H_m$ mit $N_1$ bezeichnen:
\[
\Delta_1 = \pm\,N_1\,\psi'(\varrho).
\]

Ist aber $N$ die in Bezug auf $\Omega_m$ gebildete Norm, so ist
\[
N\,\psi'(\varrho) = [N_1\,\psi'(\varrho)]^2 = \Delta_1^2,
\]
und demnach ergiebt die Formel (5) mit R\"ucksicht auf \S.\,169:
\[
\Delta = N(1-r^{-2})\,\Delta_1^2.
\]

Nach \S.\,169, (7) und (10) ist:
\weq{6}{
\begin{array}{rll}
N(1-r^{-2}) &{}=q & \text{bei ungeradem } q\\
&{}=4 & \text{f\"ur } q=2,\\[4pt]
\pm\,\Delta &{}=q\,\Delta_1^2 & \text{bei ungeradem } q\\
&{}=4\,\Delta_1^2 & \text{f\"ur } q=2.
\end{array}
}
